\section{讨论与限制}

\subsection{三条结论合起来说了什么}

三条结果指向同一个结构：本书写的信息战是**程序密集、机构隐形、晚期内耗加剧**的。
程序密集指公开渠道词（奏疏、弹劾、上疏）的量级远高于隐秘渠道的常设词；
机构隐形指厂卫系 29 人在关系网络里和平均线重合（比值 $0.99$）；
晚期内耗加剧指攻讦占比 6.4\% $\to$ 16.0\%，而对照的敌对冲突 22.7\% $\to$ 22.1\% 不动。
这三条合起来支持一个判断：通俗史著把明代的帝国治理写成了一部「以人名录展开的程序史」，
信息机构在其中只作背景板，真正承载结构的是具体的权力代理人。

\subsection{首要效度威胁：同源抽取造成的自证}

本篇的首要威胁不是样本小，而是**关系类别与派系标签都由这份语料反查导出**。
「攻讦随叙事推进上升」完全可以是抽取顺序的副产物：作者写到后半部时对人物relation更熟，
自然多标几条攻讦边；分部占比用的是带章节归属关系的分母（第 1--2 部 388 条、第 6--7 部 425 条），
分母本身也随叙事变密，虽然 2.48 这个比值对分母做了部分抵消，但不能排除共线。
能给出的缓解只有一条：年份代理口径（旁证，攻讦 11.8\% $\to$ 39.4\% $\to$ 48.8\% 单调）
与分部口径同号，两路独立刻度给出同向结果，但这仍然不是独立语料的检验。
要真正排除自证，需要第二份明史著（如《明史纪事本末》类）做平行抽取，
比较两条语料下攻讦占比的时间形状是否同构——本文不做，留作下一步。

\subsection{「信息战」是现代类比，不是史实定性}

必须重复一次，因为它最容易在转述中丢掉。明代没有近代新闻业、没有舆论场，也没有「信息战」这个题名。
本文的「信息战」只指三个可测动作的集合：信息生产（奏疏、妖书）、信息传播（谣言、流言、传单）、
信息封锁与封口（黑名单、告密、诏狱）。在这个限定下使用这个词是类比；一旦说成「明代存在 information
warfare」，就超出了本文数据的支持范围。厂卫描写在通俗史著里的文学化成分尤其要警惕：
本文所有涉及机构行为的判断都只落在本书语料内，从不外推到明代档案里的机构实况。

\subsection{抽取口径与史实的区分}

三类口径问题在本篇都出现过，都在正文标了出处：
（a）$\texttt{chapter}$ 字段里 40.4\% 是「推导」，攻讦类只有 0.5\%，
两种归属率差 39.9 个百分点，说明**抽取覆盖面本身不均匀**；
（b）派系折叠到根类会让「同阵营」虚高到 76.7\%，完整串口径才是 14.0\%；
（c）12 类关系是人造枚举，不是明代官制分类，所以「8.3\%」只能读作「本抽取口径下」。

\subsection{交付物边界}

本文交付的是**度量口径与一个具体的数据集**（`.dump/\_stats\_info\.json` 与三个作图脚本），
不是完成的信息史。要做成真正的信息史，还需要：把 1231 人的籍贯、任官年、起事的科名标准化，
把匿名文书单独开一类关系，把「监听—告发—处置」建成三元而非二元的事件结构。
这三项都不在本文的数据范围内，本文只把它们点名为缺口。

\subsection{引用完整性}

参考文献 10 条，逐条经 OpenAlex / Crossref 核验；核验中共享脚本对 monograph 全判 UNVERIFIED
（只查第一作者 + OpenAlex 429 限流，试点篇 11 条仅 1 条通过），
本文改用整条记录全作者匹配 + 退避重试的局部核验器，判据其余不变（详见 AUDIT.md）。
原书语料因无 DOI 按 BOOK\_NO\_DOI 处理，并注明版本。
